%!TEX program = lualatex
% 보편적 신앙 — classical essay layout (white page, numbered sections, footnotes)
\documentclass[11pt]{article}

\usepackage{iftex}
\ifLuaTeX\else
  \PackageError{faith}{This document requires LuaLaTeX}{Run lualatex.}
\fi

% --- 제목·이름·연도: 여기서 바꾸세요 ------------------------------------------
\newcommand{\essaytitle}{보편적 신앙}
\newcommand{\essayauthor}{길준혁}
\newcommand{\essayyear}{2026}

\usepackage{fontspec}
\usepackage{xcolor}
\usepackage{geometry}
\usepackage[tracking=true,protrusion=false,expansion=false]{microtype}
\usepackage{amsmath}
\usepackage{graphicx}
\usepackage{caption}
\usepackage{booktabs}
\usepackage{array}
\usepackage{titlesec}
\usepackage{pgfplots}
\pgfplotsset{compat=1.18}
\usetikzlibrary{arrows.meta,calc}

% Garamond carries the Latin letters; Hangul falls back to a Myeongjo face,
% with matching bold and a gently slanted italic.
\directlua{
  luaotfload.add_fallback("hangul",   {"Noto Serif CJK KR:mode=harf;"})
  luaotfload.add_fallback("hangulbf", {"Noto Serif CJK KR/B:mode=harf;"})
  luaotfload.add_fallback("hangulit", {"Noto Serif CJK KR:mode=harf;slant=0.2;"})
}
\defaultfontfeatures{Renderer=HarfBuzz}
\setmainfont{EB Garamond}[
  Numbers={Lining,Proportional},
  Ligatures={TeX,Common},
  UprightFeatures={RawFeature={fallback=hangul}},
  BoldFeatures={RawFeature={fallback=hangulbf}},
  ItalicFeatures={RawFeature={fallback=hangulit}},
  BoldItalicFeatures={RawFeature={fallback=hangulbf}},
]

% --- page ---------------------------------------------------------------------
\definecolor{paper}{HTML}{FFFFFF}
\definecolor{ink}{HTML}{111111}
\definecolor{inkgray}{HTML}{6B6B6B}
\definecolor{rulegray}{HTML}{9A9A9A}
\color{ink}

\geometry{a4paper, top=30mm, bottom=30mm, left=32mm, right=32mm}

\setlength{\parindent}{1.1em}
\setlength{\parskip}{0pt}
\setlength{\emergencystretch}{3em}
\tolerance=2500
\linespread{1.35}
\pagestyle{plain}
\raggedbottom
\frenchspacing

% --- sections -----------------------------------------------------------------
\titleformat{\section}{\normalfont\large\bfseries}{\thesection.}{0.6em}{}
\titlespacing*{\section}{0pt}{3.2ex plus 1ex minus .3ex}{1.4ex plus .3ex}

% --- figures and tables -------------------------------------------------------
\captionsetup{
  font=small, labelfont=bf, labelsep=period,
  justification=justified, singlelinecheck=true,
  margin=4mm, skip=6pt
}
\captionsetup[figure]{name={그림}}
\captionsetup[table]{name={표}, position=above, skip=6pt}
\setlength{\heavyrulewidth}{0.6pt}
\setlength{\lightrulewidth}{0.35pt}
\pgfplotsset{
  plate/.style={
    axis line style={ink,line width=.6pt},
    tick align=inside,
    tick style={ink,line width=.4pt},
    tick label style={font=\footnotesize},
    /pgf/number format/assume math mode=true,
    label style={font=\small\itshape},
    every axis plot/.append style={line cap=round},
  }
}

% PCA sample: 90 points, 10 dimensions, two dominant latent factors (numpy seed 7)
\def\pcascatter{(-0.149,-0.325) (0.139,0.652) (0.343,0.569) (-0.342,-0.818) (0.401,0.460) (-0.824,-0.485) (-0.286,0.412) (-0.032,-0.602) (1.281,0.196) (1.709,0.708) (1.947,0.177) (1.345,-0.267) (-0.290,-0.188) (2.441,0.195) (-0.274,-0.209) (1.500,0.185) (0.957,0.345) (-1.133,0.318) (-0.005,-0.692) (0.504,0.106) (-0.480,-0.108) (0.967,-0.062) (-1.660,0.729) (-1.120,-0.198) (0.652,-1.322) (-1.108,0.756) (-0.199,-0.381) (0.069,-0.539) (-0.009,-0.448) (-2.016,0.415) (-0.417,0.169) (-0.361,0.596) (0.486,0.131) (-1.016,-0.783) (1.176,0.553) (-1.138,1.104) (0.120,-0.157) (-1.531,-0.543) (0.059,0.116) (0.211,-1.064) (0.333,0.156) (-0.757,-0.122) (-0.130,0.490) (-0.377,0.106) (-1.407,-0.592) (-0.023,-0.535) (0.258,-0.924) (-0.176,-0.465) (1.239,-0.166) (1.717,0.986) (-0.081,0.292) (-0.246,-1.563) (0.644,0.346) (-0.385,-0.353) (0.080,-0.153) (-0.947,-0.619) (0.971,-0.083) (-0.310,0.588) (-0.565,0.488) (-1.224,-0.076) (-0.376,0.177) (-0.002,1.118) (0.945,-0.373) (2.326,-0.374) (1.719,-0.544) (0.763,-0.517) (-0.423,0.962) (-1.658,-1.205) (-0.107,0.030) (-0.111,0.512) (-1.412,0.136) (-0.178,0.515) (0.333,0.779) (-1.739,-0.160) (-1.301,-0.149) (0.388,0.087) (0.468,-0.090) (0.215,0.075) (1.229,0.349) (-2.129,0.303) (1.000,-0.402) (-1.846,0.906) (-0.026,0.393) (1.759,-0.490) (0.053,-0.013) (0.685,-0.316) (0.376,-0.014) (0.848,0.804) (-1.730,0.210) (-0.604,-0.209)}
\def\pcascree{(1,72.1) (2,21.4) (3,1.2) (4,1.0) (5,0.9) (6,0.9) (7,0.7) (8,0.6) (9,0.6) (10,0.6)}
\def\pcacum{(1,72.1) (2,93.5) (3,94.6) (4,95.6) (5,96.6) (6,97.5) (7,98.2) (8,98.8) (9,99.4) (10,100.0)}
\def\pcascreetwo{(1,72.1) (2,21.4)}
\def\pcatwo{93}

\usepackage[hidelinks]{hyperref}
\hypersetup{pdftitle={\essaytitle}, pdfauthor={\essayauthor}}
\urlstyle{same}


\begin{document}

% =============================================================================
\begin{center}
  {\fontsize{22}{28}\selectfont\textls[40]{\essaytitle}}\par
  \vspace{1.4em}
  {\large\essayauthor}\par
  \vspace{0.35em}
  {\normalsize\essayyear}\par
\end{center}

\vspace{2.2em}

\noindent 나는 하나님의 존재를 완전히 증명할 수 있다고 생각하지 않는다.
그렇다고 그런 하나님을 믿는 게 비합리적이라고도 생각하지 않는다. 내가 생각하는
신앙은 오히려 이런 쪽에 가깝다. 유한한 인간이 무한한 존재를 완전히 이해할 수
없다는 걸 인정하면서도, 그렇다고 그 가능성을 닫아버리지는 않고 그 안에서
살아보는 것.

이걸 설명할 때 나한테 가장 잘 맞는 비유가 수학에서의 무한이다. 우리는 모든
의미에서 무한을 직접 경험할 수 없다. 자연수를 하나씩 계속 센다고 해서 어느
순간 무한에 도착하는 것도 아니다. 끝없이 이어지는 소수점 아래 자릿수 역시
우리가 그 끝을 실제로 확인할 수는 없다.

\begin{figure}[htb]
\centering
\begin{tikzpicture}[x=1cm,y=1cm,line width=.45pt,color=ink,
                    font=\small]
  % ---- the finite: a whole that fits inside one boundary ----
  \begin{scope}[shift={(0,0)}]
    \draw[line width=.6pt] (0,0) ellipse[x radius=2.05,y radius=1.25];
    \draw[inkgray,line width=.25pt] (0,0) ellipse[x radius=1.93,y radius=1.13];
    \foreach \p/\n in {(-1.2,.35)/1,(-.55,.75)/2,(.25,.55)/3,(1.05,.7)/4,
                       (-.85,-.4)/5,(0,-.2)/6,(.85,-.25)/7,(-.15,-.8)/8}{
      \fill \p circle[radius=1.5pt];
      \node[font=\scriptsize,anchor=south west,inner sep=1pt] at \p {\n};
    }
    \node at (0,-1.75) {\textls[80]{\textsc{Finite}} \,· 유한};
    \node[font=\footnotesize\itshape,align=center] at (0,-2.3)
      {경계 하나 안에 전체가 담긴다};
  \end{scope}
  % ---- the infinite: a count that never arrives ----
  \begin{scope}[shift={(3.2,0)}]
    \draw[line width=.6pt] (0,.55) -- (6.2,.55);
    \foreach \i in {1,...,14}{
      \pgfmathsetmacro{\xx}{6.2*(1-0.84^\i)/(1-0.84^14)*0.93}
      \pgfmathsetmacro{\op}{max(0.12,1-0.065*\i)}
      \pgfmathsetmacro{\rr}{max(0.5,1.6-0.08*\i)}
      \fill[opacity=\op] (\xx,.55) circle[radius=\rr pt];
      \ifnum\i<7
        \node[font=\scriptsize,above,opacity=\op] at (\xx,.6) {\i};
      \fi
    }
    \draw[-{Stealth[length=5pt]},line width=.6pt] (6.2,.55) -- (6.9,.55)
      node[right] {$\infty$};
    \node[font=\scriptsize\itshape,inkgray] at (4.9,1.25) {끝에 닿지 않는다};
    % a decimal whose digits recede
    \node[anchor=west] at (0,-.45) {$\tfrac{1}{3} = 0.$};
    \foreach \i in {1,...,18}{
      \pgfmathsetmacro{\op}{max(0.08,1-0.055*\i)}
      \pgfmathsetmacro{\sz}{max(4.5,9.5-0.28*\i)}
      \node[anchor=west,inner sep=0,opacity=\op,
            font=\fontsize{\sz}{\sz}\selectfont]
        at ({0.98+0.33*\i-0.0055*\i*\i},-.45) {3};
    }
    \node[inkgray] at (6.55,-.45) {$\cdots$};
    \node at (3.3,-1.75) {\textls[80]{\textsc{Infinite}} \,· 무한};
    \node[font=\footnotesize\itshape,align=center] at (3.3,-2.3)
      {세어도, 써 내려가도 끝나지 않는다};
  \end{scope}
  \draw[inkgray,line width=.3pt] (2.6,-2.4) -- (2.6,1.3);
\end{tikzpicture}
\caption{유한과 무한. 왼쪽의 유한한 모임은 하나의 경계 안에 전부 들어온다.
오른쪽의 자연수와 순환소수는 아무리 이어 써도 마지막 항에 이르지 않는다.
우리는 그 끝을 본 적이 없지만, 유한한 앞부분만으로도 무한이 무엇을
뜻하는지 안다.}
\label{fig:finite}
\end{figure}

그런데도 우리는 무한이라는 개념을 이해하고 사용한다. 자연수를 계속 더해도
마지막 수가 없다는 걸 알고, 어떤 수의 소수점 아래 자릿수가 끝없이 이어진다는
것도 이해한다. 즉 무한을 직접 본 적은 없지만, 유한한 것들을 통해 무한이 무엇을
뜻하는지 어느 정도는 알 수 있다(그림~\ref{fig:finite}).
나는 하나님도 비슷하다고 생각한다.

% =============================================================================
\section{유한한 것과 무한한 것은 같은 방식으로 알 수 없다}

나는 인간의 이성이나 오감을 믿지 않는 사람이 아니다. 오히려 반대에 더 가깝다.
이성과 오감이야말로 지극히 인간적인 감각이고, 오히려 영(靈)의
세계는 나한테 낯설다. 내 앞에 책상이 있는지, 사람이 있는지, 어떤 일이 실제로
일어났는지 같은 건 인간이 충분히 판단할 수 있다. 보고, 듣고, 만지고,
논리적으로 생각하고, 다른 사람과 경험을 비교하고 검증할 수도 있다. 의리와
사랑도 여기에 포함된다. 이것들은
전부 우리와 마찬가지로 유한한 세계 안에 있는 것들이다.

그런데 무한은 조금 다르다. 여기서 “완전히 이해할 수 없다”는 말은 무한에 대해
아무것도 모른다거나 생각조차 할 수 없다는 뜻이 아니다. 오늘날 수학은 무한이라는
개념을 전혀 불편해하지 않는다. 자연수처럼 하나씩 셀 수 있는 무한(가산 무한)과
실수처럼 셀 수조차 없는 무한(비가산 무한)을 구별하고, 어떤 무한이 다른 무한보다
“더 크다”는 것까지 엄밀하게 따져 놓았다. 무한에 대한 우리의 개념적 이해는 꽤
깊다.

하지만 실제로 계산을 하는 순간 사정이 달라진다. 손으로 하든, 계산기로 하든,
컴퓨터나 \textsc{gpu}로 하든, 계산은 결국 물리적인 세계 안에서 일어나고, 물리적인
세계는 유한하다. 그래서 우리는 유한한 자릿수의 숫자로, 유한한 횟수의 연산만 할
수 있다. 컴퓨터는 $\pi$의 자릿수를 끝까지 저장할 수 없고, 보통의 계산에서는
$1/3$조차 어딘가에서 잘라낸 근삿값으로 담는다.\footnote{예를 들어 컴퓨터에서
$0.1 + 0.2$는 $0.300000000\allowbreak 00000004$가 된다. $0.1$을 이진수로 쓰면 소수점 아래
자릿수가 끝없이 이어지는데, 이를 유한한 비트에서 잘라 담기 때문이다.}
$1/3$을 ‘1과 3의 비’로, $\pi$를 하나의 기호로 들고 다니며 정확하게
계산하는 프로그램도 있다. 하지만 그것도 결국 유한한 비트로 적어 둔 이름과
규칙일 뿐, 무한한 자릿수를 펼쳐 담은 것은 아니다. 컴퓨터가 무한을
\emph{모르는} 게 아니라, 무한을 \emph{담을} 수 없는 것이다.

내가 말하는 인간의 이해도 이와 비슷하다. 나는 무한한 존재에 대해 생각할 수 있고,
그 성질을 어느 정도 짐작하고 이야기할 수도 있다. ‘하나님’이라는 이름으로 그
존재를 가리킬 수도 있다. 하지만 그 존재 전체를 내 유한한 머릿속에 남김없이
펼쳐 담을 수는 없다. 만약 그럴 수 있다면 그 존재는 애초에 무한하지 않았다는
뜻일 것이다. 그래서 하나님이 무한한 존재라면, 내가 하나님을 완전히 이해하지
못한다는 건 당연한 일이다.

이 글에서 내가 계속 현실을 강조하게 될 텐데, 이것도 같은 맥락이다.
무한을 아무리 잘 이해하더라도 실제 계산은 유한한 숫자로 해야 하듯이, 무한한
존재를 향해 열려 있더라도 내가 실제로 살아내는 건 유한한 하루와 유한한 선택들이다.

그렇다고 하나님에 대해 아무것도 알 수 없다는 뜻은 아니다. 우리가 무한에
도달하지 못하면서도 무한을 이해할 수 있듯이, 인간도 자기에게 주어진 유한한
경험을 통해 무한한 존재를 조금이나마 느끼고, 이해하고, 추론할 수 있다고
생각한다(그림~\ref{fig:asymptote}).

\begin{figure}[htb]
\centering
\begin{tikzpicture}
\begin{axis}[
  plate,
  width=11.8cm, height=6.6cm,
  xmin=-3, xmax=1.6, ymin=0, ymax=21,
  axis lines=left,
  xtick={-3,-2,-1,0,1}, ytick={0,2,4,8,16},
  xlabel={$x$}, ylabel={$f(x) = \dfrac{1}{1-x}$},
  ylabel style={rotate=-90,at={(axis description cs:-0.02,1.02)},anchor=south west},
  xlabel style={at={(axis description cs:1,-0.02)},anchor=north east},
  clip=false,
]
  % the asymptote: where f would be infinite, and which no step reaches
  \addplot[inkgray,dashed,line width=.45pt] coordinates {(1,0) (1,21)};
  \node[font=\footnotesize\itshape,anchor=west,align=left]
    at (axis cs:1.06,15.5) {$x=1$\\$f$가 무한이 되는 자리\\닿지 않는 경계};
  \addplot[ink,line width=.75pt,domain=-3:0.953,samples=260]
    {1/(1-x)};
  \draw[-{Stealth[length=5pt]},ink,line width=.55pt]
    (axis cs:0.953,21) -- (axis cs:0.957,22.4);
  % finite values, each one bounded, that together have no ceiling
  \foreach \k in {1,...,4}{
    \pgfmathsetmacro{\xk}{1-2^(-\k)}
    \pgfmathsetmacro{\yk}{2^\k}
    \addplot[rulegray,line width=.3pt,densely dotted]
      coordinates {(-3,\yk) (\xk,\yk)};
    \addplot[only marks,mark=o,mark size=1.8pt,
             mark options={ink,line width=.45pt,fill=paper}]
      coordinates {(\xk,\yk)};
  }
  \node[font=\footnotesize\itshape,anchor=east,align=right]
    at (axis cs:0.35,12) {$f(x_n) = 2^n$\\$2, 4, 8, 16, \dots$ 하나하나는 유한하지만,\\
    위로는 끝이 없다};
  \draw[inkgray,line width=.3pt,-{Stealth[length=3.5pt]}]
    (axis cs:0.37,12) -- (axis cs:0.9,15.2);
\end{axis}
\end{tikzpicture}
\caption{점근선 --- 무한을 향해 다가감. 걸음 $x_1, x_2, \dots$마다 얻는 값
$f(x_n) = 2, 4, 8, 16, \dots$은 하나하나가 모두 유한하지만 위로는 한계가 없다.
$f$가 무한이 되는 경계 $x=1$에는 결코 서지 못해도, 그 값들이 어디를 향하는지는
분명히 읽힌다. 도달하지 못한다는 것과 방향을 알 수 없다는 것은 다른 이야기다.}
\label{fig:asymptote}
\end{figure}

그래서 나는 하나님을 완전히 증명하는 걸 신앙의 출발점으로 삼고 싶지 않다. 또
천국이나 영생 같은 삶 이후의 문제를 신앙의 근거나 핵심 요인으로 끌어들이고
싶지 않다. 내가 매일 살아내는 삶도 이해하지 못하는데 어찌 삶 이후를 논하겠는가?
또 삶 이후에 어떤 금은보화가 날 기다리고 있다 한들, 그게 내 현재 삶의
인센티브가 되게 하고 싶지 않다. 마치 전쟁의 승패와 무관하게 장수의 용맹함은 그
자체로 높이 사야 마땅하고, 화재에서 구출해낸 아기가 살아야만 소방관의 희생이
고결한 게 아닌 것처럼.

때문에 내 관점은 “진실이 무엇이든 다 좋다. 그리고 그 진실이 무엇인지 우리가
부분적으로 알아냈다 치자. 하지만 어찌됐든 나는 진실을 완전히 이해하지 못하는
상태로 현실을 계속해서 살아내야 하는 존재”임을 인정함에서 시작한다. 내가
이해하지 못하는 삶 이후의 무언가를 위해 현실의 무언가를 포기하거나 희생하진
않겠다. 고집이라기보단, 정말 그러할 마땅한 이유를 모르겠다.

% =============================================================================
\section{증명하지 못한다고 해서 없는 건 아니다}

그렇다고 “어차피 증명할 수 없으니까 그냥 믿자. 그게 신앙이다.”라고 말하고 싶은
건 아니다. 내가 중요하게 보는 건 인간의 인식이 어디까지 갈 수 있느냐는 문제다.
유한한 인간이 유한한 세계를 이해하는 데는 이성과 경험이 꽤 잘 작동한다. 그런데
무한한 존재에게까지 똑같은 방식의 증명을 요구하는 게 맞는지는 모르겠다.

그래서 나는 하나님이 존재한다는 걸 완전히 증명했다고 말하지 못한다. 그렇다고
증명하지 못했으니 하나님이 없다고 말할 수도 없다고 생각한다. 쉽게 말하면
여기서 어느 정도는 포기하는 셈이다. 하지만 내가 포기하는 건 이성이 아니다.
“정말 중요한 걸 믿으려면 그 전에 무조건 완벽한 증명이 있어야 한다”는 요구를
포기하는 것이다.

나는 유한한 인간이고 오늘 자고 일어나면 내일이라는 현실을 또 살아가야 한다.
궁극적인 진리가 뭔지 완전히 확인할 때까지 내 삶을 정지시켜 놓을 수는 없다.
나는 오늘도 사람을 만나고, 사랑하거나 미워하고, 감사하거나 불평하고, 성공과
실패를 겪으면서 살아간다. 그렇다면 결국 나에게 유일하게 남는 유의미한 질문은
이것이다.
\begin{center}
\itshape “궁극적인 진실을 완전히 알 수 없는 상황에서\\
나는 어떤 방향으로 살아갈 것인가?”
\end{center}

% =============================================================================
\section{유용하다고 해서 참인 건 아니다}

나는 하나님을 믿는 게 삶에 도움이 되니까 하나님이 실제로 존재한다고 주장하고
싶은 게 아니다. 수학에서 무한이라는 개념을 갖고 극한이나 미적분 같은
도구가 발전했다고 해서, 그 사실 자체가 무한의 존재를 증명해주는 건 아니다.

하나님도 마찬가지다. 내가 하나님을 믿은 뒤에 더 감사하게 되고, 삶이 더
풍성해졌다고 해서 그게 하나님의 존재를 증명하진 않는다. “나한테 도움이
된다”와 “그게 실제로 참이다”는 다른 문제이고, 나는 그 둘을 섞고 싶지 않다.

그리고 반대도 마찬가지라고 생각한다. 어떤 믿음이 완전히 증명되지 않았다는
이유만으로 그 믿음이 내 삶에서 아무 역할도 해서는 안 된다는 결론 역시 나오지
않는다. 다시 강조하지만 결국 나는 현실을 살아야 한다. 그리고 궁극적인 진실에
대해 완전한 이해나 확신을 가질 수 없는 상황이라면, 내가 어떤 세계관을 가지고
사는 게 실제 내 삶을 더 풍성하게 만드는지는 충분히 고려할 만한 문제라고
생각한다.

% =============================================================================
\section{설령 그게 어느 정도 낭만이라 해도}

나는 궁극적인 진실과 우리가 실제로 살아가는 현실 사이에는 어느 정도 차이가 있을
수 있다고 생각한다. 만약 어떤 형이상학적 진실이 정말 존재하는데, 그게 유한한
인간의 삶과 완전히 동떨어져 있고, 우리가 알 수도 경험할 수도 없고, 실제 삶에
아무 영향도 주지 않는다면, 굳이 그 진실을 우리가 추구해야 하나 싶다.

반대로 어떤 믿음이 궁극적으로는 일종의 낭만이나 환상일 가능성을 완전히 배제할
수 없더라도, 그 믿음이 사람으로 하여금 더 감사하고, 더 사랑하고, 자기 삶을 더
충실하게 살아가게 만든다면 단지 “증명되지 않았다”는 이유만으로 꼭 버려야 할까?
나는 그렇지는 않다고 생각한다.

물론 이건 내가 거짓이라고 아는 걸 일부러 믿겠다는 말은 아니다. 나는 그게
거짓인지 참인지 최종적으로 모르는 상태에 있다. 그 상태에서 내가 할 수 있는 건
현재 나한테 주어진 이성과 경험을 가지고 가장 합리적이라고 생각되는 방향, 그리고
실제 삶을 가장 풍성하게 만드는 방향을 선택하는 것이다.

그리고 나중에 더 합리적인 설명이나 더 나은 길이 나온다면 내 생각을 바꿀 의향도
있다. 신앙 때문에 생각을 멈추고 싶지는 않다. 예수님을 믿는다고 고백함으로
구원을 얻고 내 삶과 그 이후의 문제가 전부 해결됐다고 여기고 싶지도 않다.

% =============================================================================
\section{그런데 왜 ‘더 나은 삶’이 진실과 관계 있다고 생각하는가}

사실 내 주장은 전부 하나의 전제 위에 서 있다. 우리를
만든 창조주가 있다는 전제다. 그리고 그 창조주는 앞에서 내가 증명할 수 없다고 말한
바로 그 대상이다. 그러니 이 장은 “창조주가 있다면”으로 시작하는
조건문으로 읽혀야 한다.

이 전제를 왜 붙드느냐고 묻는다면 “믿었더니 삶이 풍성해져서”라고 답하고 싶지는
않다. 그렇게 답하면 삶이 풍성하다는 이유로 창조주를 믿고, 창조주가 있다는 이유로
풍성한 삶이 진실과 닿아 있다고 말하는 순환논리가 된다. 삶의 풍성함은
내가 이 전제를 \emph{택하는} 이유는 될 수 있어도, 이 전제가 \emph{참이라는}
근거는 되지 못한다. 내가 이 전제를 붙드는 건 증명도 반증도 닿지 않는 위치에서,
유한한 경험들이 가리키는 쪽으로 한 걸음을 내디딘 것에 가깝다. 이 전제
위에서라면, 나는 궁극적인 진실과 궁극적인 선이 완전히 반대 방향에 있을
거라고 보지 않는다.

만약 우리를 만든 창조주가 실제로 있고, 우리가 단순히 아무 의미 없이 생겨난
존재가 아니라 그 창조주의 피조물이라면, 창조에는 어떤 의도가 있을 거라고
생각한다. 물론 내가 그 의도를 완전히 알 수 없다. 유한한 피조물이 무한한
창조주의 의도를 다 이해하는 건 불가능하다. 하지만 전혀 알 수 없는 것과 조금이라도
알아차릴 수 있는 건 다르다.

우리는 살아가면서 사랑, 은혜, 양심, 관계, 감사, 창조성 같은 것들을 경험한다.
나는 이런 것들을 통해서 무한한 창조주의 의도가 어떤 방향에 있는지 유한한 우리가
아주 조금이라도 짐작할 수 있다고 생각한다.

피조물에게 좋은 삶이라는 건 결국 자기를 만든 의도와 점점 멀어지는 삶보다는, 그
의도를 조금이라도 알아채고 그 방향으로 살아가려고 노력하는 삶일 거라고
생각한다. 아무리 무한이 멀리 떨어져 있어도 방향만 알면 그곳을 향해 유한한
발걸음으로 나아갈 수 있다. 내가 “궁극적인 진실은 결국 우리의 삶을 더 나은
방향으로 이끌 것”이라고 생각하는 이유도 여기에 있다.

\begin{figure}[htb]
\centering
\begin{tikzpicture}
\begin{axis}[
  plate,
  width=0.94\linewidth, height=7.4cm,
  xmin=0, xmax=100, ymin=0, ymax=10.2,
  axis lines=left,
  xtick={0,20,...,100}, ytick={0,2,...,10},
  xlabel={걸음 $n$}, ylabel={$a_n$},
  xlabel style={at={(axis description cs:1,-0.09)},anchor=north east},
  ylabel style={rotate=-90,at={(axis description cs:-0.01,1.0)},anchor=south east},
  grid=major,
  major grid style={rulegray!40,line width=.2pt},
  clip=false,
]
  % the finite ceiling
  \addplot[inkgray,dashed,line width=.45pt] coordinates {(0,6) (100,6)};
  \node[font=\footnotesize\itshape,anchor=south east,inkgray]
    at (axis cs:100,6.05) {유한한 상한 $L = 6$};
  % converging: fast, impressive, and finished
  \addplot[ink,line width=.7pt,densely dashed,domain=0:100,samples=160]
    {6*(1-exp(-x/6))};
  % diverging: slow, unimpressive, and unbounded
  \addplot[ink,line width=.85pt,domain=0:100,samples=200]
    {0.8*sqrt(x)};
  \draw[-{Stealth[length=5pt]},ink,line width=.7pt]
    (axis cs:99,7.96) -- (axis cs:104,8.16);
  % crossing
  \addplot[only marks,mark=o,mark size=2.2pt,
           mark options={ink,line width=.5pt,fill=paper}]
    coordinates {(56.3,6.0)};
  \node[font=\footnotesize\itshape,anchor=north west,align=left]
    at (axis cs:57,5.7) {$n \approx 56$:\\발산하는 쪽이 상한을 넘는다};
  % labels
  \node[font=\footnotesize,anchor=south west,align=left]
    at (axis cs:6,6.7) {수렴 \quad $a_n = 6\,(1-e^{-n/6}) \;\to\; 6$\\
    \itshape 처음엔 훨씬 높지만, 거기서 멈춘다};
  \node[font=\footnotesize,anchor=north west,align=left]
    at (axis cs:22,2.9) {발산 \quad $b_n = 0.8\sqrt{n} \;\to\; \infty$\\
    \itshape 느리고 낮게 출발해도, 끝이 없다};
  \node[font=\small,anchor=west] at (axis cs:105,8.25) {$\infty$};
\end{axis}
\end{tikzpicture}
\caption{수렴과 발산. 어느 한 순간의 높이만 비교하면 수렴하는 수열이 앞서
보일 수 있다. 그러나 한쪽은 아무리 커도 유한한 값 $L$ 아래에 갇히고, 다른
한쪽은 무한을 향해 열려 있다. 둘은 결과적으로 전혀 다른 모양이다.}
\label{fig:diverge}
\end{figure}

절대적인 무한 앞에서 상대적인 위치는 의미 없지만 그 무한을 향해 발산하고
있다는 사실은 유의미하다(그림~\ref{fig:diverge}). 아무리
크더라도 특정 유한한 값으로 수렴하는 것과 무한대로 발산하는 건 결과적으로 아예
다른 모양이다.

창조주가 있고 우리가 그의 피조물이라면, 그 창조의 의도에 가까워지는 삶이
궁극적으로 우리를 망가뜨리는 방향일 거라고 생각하기는 어렵다. 오히려 피조물이
자기가 왜 존재하는지에 가까워질수록 더 자기답고 충만한 삶을 살게 되는 게
자연스럽다고 생각한다.

물론 이 전제가 틀렸을 수도 있다. 창조주가 없다면 이 장에서 말한 진실과 선의 연결도
함께 무너진다. 그때 내게 남는 건 앞 장에서 말한 것뿐이다. 증명되지 않은
믿음이었지만, 그 덕분에 조금 더 감사하고 조금 더 사랑하며 살았다는 것. 나는
그것만으로도 이 선택을 후회하지 않을 것 같다. 다만 그것을 이 전제가 옳았다는
증거로 삼지는 않겠다.

% =============================================================================
\section{파스칼의 내기, 판돈을 바꾸면}

여기까지 오면 떠오르는 논증이 하나 있다. 17세기의 수학자이자 철학자였던 블레즈
파스칼의 ‘내기’다.\footnote{파스칼, 『팡세』, 브룅슈비크 판 단편 233.} 파스칼은
이성만으로는 신이 있는지 없는지 판가름할 수 없다고 보았다. 그런데도 우리는 어느
한쪽에 걸 수밖에 없다. 살아간다는 것 자체가 이미 내기판에 올라와 있는 것이기
때문이다. 그렇다면 어느 쪽에 거는 게 합리적인가. 파스칼의 답은 이렇다. 신이
있다는 쪽에 걸어서 이기면 무한히 행복한 무한한 생명을 얻고, 지더라도 잃는 것은
유한하다. 반대쪽에 걸었다가 지면 무한한 것을 잃는다. 그래서 그는 “이기면 모든
것을 얻고, 져도 아무것도 잃지 않는다”고 말한다(표~\ref{tab:wager}의 (a)).

\begin{table}[htb]
\centering\footnotesize
\caption{두 개의 내기. (a)는 파스칼의 표로, 거는 대상은 믿음이고 판돈은 삶 이후의
무한이다. (b)는 이 글의 변형으로, 거는 대상은 삶의 방향이고 네 칸에는 모두 지금
이 삶에서 일어나는 유한한 것들이 적힌다.}
\label{tab:wager}
\renewcommand{\arraystretch}{1.45}
\begin{tabular}{@{}>{\raggedright\arraybackslash}p{2.8cm}
                >{\centering\arraybackslash}p{4.6cm}
                >{\centering\arraybackslash}p{4.6cm}@{}}
  \multicolumn{3}{@{}l@{}}{\small (a) 파스칼의 내기} \\[0.2em]
  \toprule
  & 신이 있다 & 신이 없다 \\
  \midrule
  믿는다      & $+\infty$\,· 무한한 생명 & 유한한 손실 \\
  믿지 않는다 & $-\infty$\,· 무한한 상실 & 유한한 이득 \\
  \bottomrule
  \addlinespace[1.4em]
  \multicolumn{3}{@{}l@{}}{\small (b) 이 글의 변형} \\[0.2em]
  \toprule
  & 창조주가 있다 & 창조주가 없다 \\
  \midrule
  그 방향으로 산다 & 진실과 선이 한 방향에 놓인다 & 조금 더 감사하고 사랑한 삶 \\
  그러지 않는다   & 존재의 이유에서 멀어진다     & 이 표로는 판가름할 수 없다 \\
  \bottomrule
\end{tabular}
\end{table}

이 논증은 내 생각과 닮은 데가 많다. 증명할 수 없다는 데서 출발한다는 점, 판단을
미룬 채로는 살 수 없다는 점, 그리고 그 상태에서 어느 쪽을 고르는 게 합리적인지
묻는다는 점이 그렇다. 앞에서 “궁극적인 진실을 완전히 알 수 없는 상황에서 나는
어떤 방향으로 살아갈 것인가?”라고 물었던 것도 결국 같은 질문이다. 다만 나는 이
내기를 그대로 받아들이지 않고 두 군데를 바꾼다.

첫째, 판돈이 다르다. 파스칼의 표에서 승부를 가르는 건 삶 이후의 무한한 보상과
상실이다. 앞에서 말했듯 나는 삶 이후의 금은보화를 지금 삶의 인센티브로 삼고 싶지
않다. 그래서 내 표의 칸에는 무한대 대신 지금 이 삶에서 일어나는 유한한 것들을
적는다.

둘째, 거는 대상이 다르다. 파스칼은 믿음에 걸라고 한다. 하지만 뒤에서 다시
말하겠지만, 어떤 명제가 믿어지는 건 내가 고를 수 있는 일이 아니다. 파스칼도 이
어려움을 알았던 것 같다. 그는 믿음이 생기지 않는 사람에게, 이미 믿는 사람들처럼
성수를 받고 미사를 드리며 살아보라고 권한다. 나는 그 권유를 조금 다르게 읽는다.
그렇게 살다 보면 믿음이 생긴다는 데 방점을 찍기보다, 어떻게 살지는 내가 고를 수
있다는 데 방점을 찍는다. 그래서 내가 거는 건 믿음이 아니라 방향이다.

이렇게 바꾸면 표는 표~\ref{tab:wager}의 (b)처럼 된다. 창조주가 있다면, 그 방향으로 사는 삶은 창조의
의도와 같은 쪽을 향한다. 앞 장에서 말한 대로 진실과 선이 거기서 만난다. 창조주가
없다면, 그 삶은 증명되지 않은 믿음 위에서 조금 더 감사하고 조금 더 사랑하며 산
삶으로 남는다. 그 방향을 택하지 않았는데 창조주가 있다면, 나는 내가 존재하는
이유에서 조금씩 멀어진 셈이다. 다만 이 손실도 무한한 형벌로 적지 않고, 지금 이
삶에서 놓친 것으로 적는다. 창조주도 없고 그 방향도 택하지 않았다면, 이 표는 그
삶에 대해 아무것도 판가름하지 못한다. 그런 사람은 다른 곳에서 자기 방향을 찾을
것이고, 그것을 내가 낮춰 볼 이유는 없다.

이 변형에서 흥미로운 건 무한이 사라지지 않고 자리를 옮긴다는 점이다. 파스칼에게
무한은 칸 안에 적힌 보상이었다. 나에게 무한은 보상이 아니라 방향이다. 발산하는
수열이 어느 한 순간의 높이가 아니라 향하는 곳으로 구별되듯이(그림~\ref{fig:diverge}),
내가 거는 것도 나중에 받을 무언가가 아니라 매일의 걸음이 향하는 쪽이다.

계산하는 방식도 달라진다. 파스칼의 내기는 흔히 기대값으로 설명된다.\footnote{우월
논증과 기대값 논증의 구분은 A.~Hájek, “Pascal's Wager,” \emph{Stanford
Encyclopedia of Philosophy} 참조.} 신이 있을
확률이 아무리 작아도 무한한 보상을 곱하면 무한이 되니, 거기에 거는 게 낫다는
식이다. 내 표에는 무한이 없으니 이 계산은 쓸 수 없다. 대신 더 단순한 구조가
남는다. 윗줄이 두 칸 모두에서 아랫줄보다 못하지 않다면, 창조주가 있을 확률을 몇으로
잡든 윗줄을 고르는 게 합리적이다. 결정이론에서 ‘우월’이라고 부르는 구조이고,
파스칼의 논증에도 이 형태가 들어 있다.

물론 여기에는 조건이 하나 붙는다. 창조주가 없는 쪽에서도 그 방향으로 사는 삶이
그렇지 않은 삶보다 못하지 않아야 한다. 파스칼이 “져도 아무것도 잃지 않는다”고
말한 자리가 바로 여기다. 나는 이 칸을 증명으로 채우지 못한다. 지금까지 살아보고
얻은 내 판단으로 채울 뿐이다. 다른 경험을 한 사람은 이 칸을 다르게 채울 수 있고,
그러면 표의 결론도 달라진다. 이 표가 무엇보다 나에게 유효한 표인 이유다.

파스칼의 내기가 흔히 받는 반론 가운데 하나는 “어느 신에게 걸 것인가”다. 걸 수
있는 신이 여럿이면 표가 두 줄로는 모자라고, 계산도 무너진다. 내 변형은 이 반론에서
조금은 자유롭다. 판돈이 올바른 신을 맞히는 데 있지 않고 삶의 방향에 있기 때문이다.
다른 전통 안에서 자란 누군가도 자기 언어로 같은 표를 그릴 수 있다.

% =============================================================================
\section{그래서 하나님을 믿는다는 것}

그래서 내게 하나님을 믿는다는 건 단순히 “하늘 어딘가에 초월적인 존재가 하나
있다”는 명제에 동의하는 것만은 아니다. 또 전통 기독교에서 말하는 “성경적
삼위일체 하나님을 믿는다”와도 차이가 있을 수 있다.

나는 유한한 존재이고, 내 이해를 넘어서는 무한한 창조주가 존재할 가능성을
받아들이는 것이다. 그리고 그 창조가 아무 의미 없는 일이 아니라 어떤 의도를
가지고 있으며, 나 역시 그 의도 안에서 존재하는 피조물이라고 받아들이는 것이다.

하나님의 은혜와 사랑을 경험한다는 것도 나는 이 안에서 이해한다. 내 삶과
사람들과 내가 가진 능력과 오늘 하루를 그냥 원래 내 것이었던 것처럼 보는 대신,
어떤 의미에서는 주어진 것으로 받아들이면 삶을 바라보는 방식이 달라진다. 감사할
이유가 많아진다. 내가 모든 걸 내 힘으로 만들어내고 소유해야 하는 존재가 아니라,
이미 많은 것을 받은 존재라는 생각도 하게 된다.

그리고 내가 창조주와 아무 상관 없는, 수많은 원자들로 이루어진 고립된 개인이
아니라 그 관계 안에 있는 존재라고 생각하면 삶은 단순히 먹고 살고 성공하고 죽는
것 이상이 된다. 창조주와 동행하는 삶을 살게 된다.

나는 그런 의미에서 하나님의 은혜와 사랑 안에서 살아가는 게 현실에서 천국을
살아가는 한 방식이라고 생각한다. 현실을 떠나는 게 아니라 같은 현실을 조금 다른
지평에서 살아가는 것이다.

여기까지는 어느 종교의 언어로도 할 수 있는 이야기다. 그렇다면 나는 왜 하필
기독교의 언어로 이 이야기를 하는가. 먼저 인정할 것은 내가 기독교 안에서
자랐다는 사실이다. 나는 한 번뿐인 삶을 이 버젼으로 살아왔고, 그만큼 다른
종교와 교리를 깊이 겪어보지 못했다. 그래서 내 삶의 경험들은 기독교의
언어로밖에 이어지지 않는다. 다른 길을 모르거나 밀어내는 게 아니다. 서로 다른
깊이로 알고 있는 것들을 나란히 놓고 우열을 가리는 게 오히려 합리적이지 않다는
것이다.

다만 그 안에서 자랐다는 이유만으로 기독교를 받아들이고 싶지는 않았다. 이 글은
주어진 환경 안에서 무엇을, 어떤 기준으로 받아들일지에 대한 기록이다. 그래서 이
글이 보편적이라면, 그건 내용이 아니라 기준의 차원에서다. 힌두교 안에서 나와
비슷하게 살아온 누군가가 있다면, 그 사람도 같은 기준으로 자기 종교를 바라볼 수
있을 것이다.

% =============================================================================
\section{신앙은 확실성보다 신뢰에 가깝다}

그래서 내게 신앙은 “나는 하나님이 존재한다는 걸 완벽하게 안다”고 말하는 데 있지
않다. 모르는 건 모른다고 인정하고 싶다. 내가 죽을 때까지 모르는 것도 있을
것이다.

나는 아직 “예수의 부활”, “구원의 확신”, “천사”, “사탄”, “영혼”, “지옥” 등을
어떻게 받아들일지 잘 모르겠다. 나는 이것들이 시대적, 종교적, 문학적, 추상적
개념이나 도구에 더 가까울 수도 있다는 가능성을 남겨두고 싶다. 이것들은 내가
서 있는 유한한 자리에서는 확인할 길이 없는 이야기들이니까.

그런데 그렇다고 신앙이 불가능한 건 아니라고 생각한다. 오히려 신앙은 내가 모르는
걸 억지로 안다고 우기는 게 아니라, 내가 어디까지 알 수 있는지를 인정하고 그
상태에서 어떤 방향으로 살아갈지를 선택하는 것에 더 가깝다.

그리고 지금까지 내가 보고 경험하고 생각한 범위에서는, 창조주와 그 의도를
인정하고 하나님의 은혜와 사랑 안에서 살아가는 게 그걸 부정하고 사는 것보다 내
삶을 더 풍성하게 만든다고 생각한다. 이건 그 전제가 참이라는 증거가 아니라,
참인지 모르는 상태에서 내가 이쪽을 고르는 이유다.

예수님의 기적, 보혈, 죽음과 부활을 부정한다는 게 아니다. 다만 나는 ‘믿는다’는
말을 조금 다르게 본다. 어떤 주장이 믿어지는 건 내가 원한다고 되는 일이 아니다.
선언하고 고백한다고 믿어지는 것도 아니고, 누군가에게 이미 믿어진 것을 내가
설득해서 믿어지지 않게 할 수도 없다. 믿음은 결심이라기보다 받아들여지는 것에
가깝고, 무엇이 받아들여지는지는 저마다의 경험과 살아가는 시대, 그리고 해석에
따라 다를 수 있다.

반면 내가 고를 수 있는 건 방향이다. 내게 예수님을 믿는다는 건 유한한 삶을
이렇게 살아가겠다는 방향을 정하는 일이다. 그렇다면 물어야 할 건 부활이 내게
믿어지느냐가 아니라, 그것이 내가 정한 방향과 어떤 관계에 있느냐다. 우리는
무한의 끝을 확인하지 않은 채 무한을 받아들이고,\footnote{집합론에서도
무한집합의 존재는 증명해서 얻는 정리가 아니라, ‘무한 공리’로 받아들이고
시작하는 전제다.} 그 위에서 극한을 구하고 미적분을 푼다. 그 끝을 확인하지
못했다고 계산이 멈추지는 않는다. 마찬가지로 내가 따르려는 삶의 방향은, 부활이
내게 믿어지든 아직 믿어지지 않든 그대로 남는다.

“그리스도께서 살아나지 않으셨다면, 우리의 선포도 헛되고, 여러분의 믿음도
헛될 것입니다”라는 바울의 말이 얼마나 무거운지 모르지 않는다.\footnote{고린도전서
15:14(새번역). 같은 뜻이 15:17에서도 반복된다.} 그래도 성경의 모든 주장을
하나하나 판가름한 뒤에야 따를 수 있다면, 나는 끝나지 않는 논쟁에 갇혀 정작
따르는 일은 시작하지 못할 것이다. 천지창조가 며칠 동안
이뤄졌는지, 홍해가 정말 갈라졌는지, 물이 정말 포도주가 되었는지가 내게
믿어지든 아니든, 그것이 그 방향을 바꾸지는 않는다. 그래서 설령 부활이 역사적
사실이 아니라 해도, 내가 예수님의 제자로 살아가지 않을 이유는 아직 뚜렷하지
않다.

예수님은 율법학자가 “가장 으뜸되는 계명이 무엇인지” 물었을 때, 첫째로 하나님을
사랑하는 것, 둘째로 네 이웃을 네 자신과 같이 사랑하는 것, 그리고 이보다 더 큰
계명은 없다고 확고하게 말씀하셨다.\footnote{마가복음 12:28--31; 마태복음
22:36--40.}

다른 기독교의 교리와 말씀을 내가 전부 이해하고 실천하지 못하더라도, 적어도 이 두
가지를 삶의 가장 중요한 방향으로 삼는다면 창조주가 의도한 방향으로 살아가는 것이
아닐까? 마치 복잡한 고차원의 데이터를 이해하기 위해 \textsc{pca}(principal
component analysis)에서 가장 중요한 두 개의 주성분 위에 데이터를 펼쳐놓고 전체적인
구조와 방향을 바라보는 것처럼 말이다(그림~\ref{fig:pca}).

\begin{figure}[htb]
\centering
\begin{tikzpicture}
\begin{axis}[
  plate,
  name=proj,
  width=8.9cm, height=5.6cm,
  xmin=-2.7, xmax=3.1, ymin=-2.1, ymax=2.1,
  axis lines=middle,
  axis line style={ink,line width=.5pt,-{Stealth[length=4pt]}},
  xtick=\empty, ytick=\empty,
  clip=false,
]
  \addplot[only marks,mark=*,mark size=1.05pt,
           mark options={inkgray,draw=inkgray}] coordinates {\pcascatter};
  % one life, read on the two principal axes
  \addplot[ink,line width=.7pt,smooth,tension=.55,
           -{Stealth[length=5.5pt]}]
    coordinates {(-1.9,-1.5) (-1.05,-1.2) (-.55,-.35) (.35,-.05) (.9,.75) (1.9,1.25)};
  \addplot[only marks,mark=o,mark size=2pt,
           mark options={ink,line width=.5pt,fill=paper}]
    coordinates {(-1.9,-1.5) (-1.05,-1.2) (-.55,-.35) (.35,-.05) (.9,.75)};
  \node[font=\footnotesize\itshape,anchor=north west] at (axis cs:-1.85,-1.55)
    {나의 궤적};
  \node[font=\footnotesize,anchor=north east,align=right,fill=paper,inner sep=1.5pt]
    at (axis cs:3.1,-.12) {PC\,1\\\itshape 하나님을 사랑하라};
  \node[font=\footnotesize,anchor=south west,align=left]
    at (axis cs:.06,2.1) {PC\,2 \quad\itshape 네 이웃을 사랑하라};
\end{axis}
\begin{axis}[
  plate,
  at={(proj.east)}, anchor=west, xshift=2.1cm,
  width=5.6cm, height=5.4cm,
  bar width=6pt,
  ymin=0, ymax=100, xmin=0.3, xmax=10.7,
  axis lines=left,
  xtick={1,...,10}, ytick={0,25,50,75,100},
  xlabel={주성분}, ylabel={설명된 분산 (\%)},
  ylabel style={font=\footnotesize\itshape},
  xlabel style={font=\footnotesize\itshape},
  tick label style={font=\scriptsize},
  clip=false,
]
  \addplot[ybar,bar shift=0pt,fill=inkgray!35,draw=ink,line width=.35pt]
    coordinates {\pcascree};
  \addplot[ybar,bar shift=0pt,fill=ink,draw=ink,line width=.35pt]
    coordinates {\pcascreetwo};
  \addplot[ink,line width=.5pt,mark=o,mark size=1.5pt,
           mark options={fill=paper}] coordinates {\pcacum};
  \node[font=\scriptsize\itshape,anchor=west,align=left]
    at (axis cs:2.9,66) {처음 두 축이\\\pcatwo\%를 설명 (예시)};
\end{axis}
\end{tikzpicture}
\caption{주성분 분석(예시 자료). 회색 점은 삶에서 내리는 유한한 선택들이고, 각
선택은 열 가지 차원을 가진다고 하자. 왼쪽은 선택들을 가장 큰 두 주성분 위에 펼친
모습, 오른쪽은 각 성분이 설명하는 분산이다. 나머지 여덟 차원은 보이지 않아도,
흰 점을 이은 나의 궤적이 어디에 서 있고 어디로 가는지는 두 축만으로 읽을 수 있다.}
\label{fig:pca}
\end{figure}

모든 차원을 다 보고 있는 것은 아니지만, 가장 중요한 두 축을 통해 적어도 내가 어느
방향에 서 있고 어디를 향해 가고 있는지는 볼 수 있다. 그래서 나는 지금은 그 방향을
선택한다.

그렇다고 그 선택 때문에 새로운 증거나 더 나은 설명을 거부하고 싶지는 않다. 더
합리적인 길이 나온다면 들을 생각도 있고, 필요하다면 내 생각을 바꿀 수도 있다.
다만 현재까지는 무한을 굳이 부정해야 할 충분한 이유를 찾지 못했다.

유한한 인간으로서 무한을 향해 열려 있고, 나를 만든 창조주의 의도를 조금이라도
알아가면서 그 방향으로 살아가려고 노력하는 게 지금 내가 할 수 있는 가장
합리적이고 인간적인 선택이라고 생각한다.

\vspace{0.6em}

\begin{center}
우린 결코 무한을 완전히 이해할 수는 없다.\\
그렇지만 그 방향을 보고 살아갈 수는 있다.\\
나는 신앙이 결국 그런 게 아닐까 생각한다.
\end{center}

% =============================================================================
\titlespacing*{\section}{0pt}{2.4ex plus .6ex minus .3ex}{0.9ex plus .2ex}
\section*{덧붙이는 말}

{\small
\noindent 이 글은 누군가를 설득하려고 쓴 글이 아니다. 믿지 않는 사람에게 기독교를
전하려는 것도, 기독교인에게 반론을 제기하려는 것도 아니다. 이 글의 독자는
무엇보다 나 자신이다. 나는 왜 지금의 신앙을 갖게 되었는지 스스로 돌아보고
싶었고, 그걸 내가 할 수 있는 만큼 논리적으로 적어보고 싶었다. 적어 내려가는
동안 나는 더 다양한 생각을 해봤고, 그 과정 자체를 즐겼던 것 같다. 그러니 혹시
이 글을 읽고 공감이 되었다면, 그건 기독교나 어떤 철학이 옳아서가 아닐 것이다.
그저 나라는 한 사람의 생각에 공감한 것뿐일 것이다.\par}

% =============================================================================
\section*{참고문헌}

{\footnotesize
\begin{list}{}{\setlength{\leftmargin}{1.6em}\setlength{\itemindent}{-1.6em}
               \setlength{\itemsep}{0.15em}\setlength{\parsep}{0pt}}
  \item 파스칼, 블레즈. 『팡세』. 1670. 브룅슈비크 판 단편 233.
  \item Hájek, Alan. “Pascal's Wager.” \emph{Stanford Encyclopedia of
        Philosophy}. \url{https://plato.stanford.edu/entries/pascal-wager/}.
  \item 『성경전서 새번역』. 대한성서공회.
\end{list}\par}

\end{document}